\section{Model, noise laws, and output contracts}\label{sec:model}

This section fixes the input model, the perturbation laws, the meaning of
an exact smoothed algorithm, and the output formats. Later sections state
their theorems in these terms. The section ends with two elementary facts
that apply to every result: a certified bound for the original objective,
and an example showing why a finite law forces an exact fallback.

\subsection{Instances and input length}\label{sec:model:input}

All data are rational and encoded in binary. A \emph{base instance}
consists of an objective $F_0$, a feasible set $X\subseteq\R^n$, a noise
half-width $\sigma>0$, and the structural data required by the relevant
theorem: a quadratic convexifier, a tree decomposition, a core, curvature
bounds, or certificates. Its \emph{base length} $I$ is the total binary
length of these data. It is measured before any noise is drawn; the bit
lengths of sampled coefficients are accounted for separately.

\paragraph{Objectives.}
A quadratic objective is $F_0(x)=\frac12x^{\mathsf T}Ax+b^{\mathsf T}x+c$
with a symmetric rational matrix $A$. A polynomial objective of total degree
at most $d$ is given by an explicit list of monomials, that is, rational
coefficients with exponent vectors, or as a sum of factors, each such a list
on a stated set of variables (its \emph{scope}). In every polynomial theorem
$d$ is a fixed constant, and $\poly_d(\cdot)$ denotes a polynomial whose
exponent depends on $d$ only. Objectives are never given as arithmetic
circuits. Circuits can be exponentially more compact than monomial lists,
and they change the complexity of evaluation and of exact output;
\cref{sec:lim:points} uses this difference.

\paragraph{Domains.}
A \emph{mixed box} is $X=\prod_{i=1}^nX_i$, where each $X_i$ is either a
continuous interval $[\ell_i,u_i]$ or a \emph{native integer} interval
$[\ell_i,u_i]\cap\Z$. Native integer bounds are encoded in binary, so the
number of labels $u_i-\ell_i+1$ can be exponential in $I$. Integer bounds
are rounded inward, empty domains are rejected, and fixed coordinates are
substituted before any algorithm starts. After this preprocessing every
remaining coordinate has width $w_i=u_i-\ell_i>0$, and the
\emph{continuous hull} of $X$ is $\prod_i[\ell_i,u_i]$. We write $n_c$ and
$n_z$ for the numbers of continuous and native integer coordinates. A
\emph{mixed polytope} is $P\cap(\Z^{n_z}\times\R^{n_c})$ for a bounded
rational polyhedron $P$ given by linear inequalities. Products of simplices,
order polytopes, polynomial graph constraints, integer flows, and totally
unimodular systems are defined where they are used.

\paragraph{Structural and numerical parameters.}
The structural parameters are the negative inertia $k=n_-(A)$ of a quadratic
objective, the bag size $p$ (treewidth plus one) of a supplied tree
decomposition whose bags contain all factor scopes, the size $k$ of a
supplied continuous core, and, in the integer theorems, the integer
dimension. The numerical parameters are curvature and width quantities
compared with $\sigma$. Examples are an upper bound $L$ on coordinate second
derivatives, $\nu=\max\{0,-\lambda_{\min}(A)\}$, a convexifier scale
$\alpha$, and coordinate widths $w_i$. Each theorem states on which set a
curvature bound must hold. A numerical parameter enters a bound through its
magnitude, not through its encoding length.

\paragraph{Status of premises.}
Each premise of a theorem has one of three statuses.
It is \emph{verified} if the algorithm decides it within the stated work
bound; examples are exact definiteness tests of rational matrices and a
curvature bound computed from monomial coefficients. It is \emph{certified}
if the input contains a proof whose verification cost is included in $I$ and
in the stated bound. It is \emph{promised} if correctness relies on it but
the algorithm need not detect a false premise. A premise that is
promised for correctness may also be needed only for the running-time
analysis; theorems say which.

\subsection{Linear perturbations and finite laws}\label{sec:model:noise}

Every result perturbs only the linear part of the objective. The
\emph{sampled objective} is
\begin{equation}\label{eq:model:sampled}
 F_\gamma(x)=F_0(x)+\gamma^{\mathsf T}x,
\end{equation}
for a random vector $\gamma\in\Q^n$ drawn once. The feasible set, the
nonlinear coefficients, and all structural data are not perturbed.

\begin{definition}[Grid law]\label{def:model:grid}
For a rational $\sigma>0$ and a power of two $M\ge2$, let
\[
 \Lambda_{\sigma,M}=\Bigl\{-\sigma+\frac{2\sigma j}{M-1}:j=0,1,\ldots,M-1\Bigr\}
\]
be the $M$ equally spaced points of $[-\sigma,\sigma]$, including both
endpoints, and let $U_{\sigma,M}$ be the uniform distribution on
$\Lambda_{\sigma,M}$.
\end{definition}

A draw from $U_{\sigma,M}$ uses exactly $\log_2M$ fair random bits, and every
atom has binary length $O(\log M+\operatorname{size}(\sigma))$. The atoms
are equally spaced at distance $2\sigma/(M-1)$, so a closed interval of
length $\lambda$ contains at most $\lambda(M-1)/(2\sigma)+1$ of them. Hence
\begin{equation}\label{eq:model:interval}
 U_{\sigma,M}(J)\le\frac{\lambda}{2\sigma}+\frac1M
 \qquad\text{for every closed interval $J$ of length $\lambda$.}
\end{equation}
The first term is the bound for the continuous uniform law on
$[-\sigma,\sigma]$. The second is an \emph{atomic term}; it does not
decrease with $\lambda$, and every finite-law estimate in this paper carries
a term of this kind.

\begin{definition}[Perturbation models]\label{def:model:perturbation}
Let $\mathcal D$ be a finite rational law on $\Q$. The following models
specify $\gamma$ in \eqref{eq:model:sampled}.
\begin{enumerate}[label=(\roman*)]
\item \emph{Ambient}: $\gamma_1,\ldots,\gamma_n$ are independent with law
$\mathcal D$.
\item \emph{Core-only}: for a supplied core $C\subseteq[n]$, the coefficients
$\gamma_i$, $i\in C$, are independent with law $\mathcal D$, and
$\gamma_i=0$ for $i\notin C$.
\item \emph{Aligned}: for a supplied rational matrix $T\in\Q^{k\times n}$,
$\gamma=T^{\mathsf T}\xi$, where $\xi_1,\ldots,\xi_k$ are independent with
law $\mathcal D$.
\end{enumerate}
\end{definition}

The models are not nested. Aligned coefficients are generally dependent in
the original coordinates and are supported on the row space of $T$; an
aligned theorem says nothing about independent ambient coefficients, and
conversely. Core-only noise leaves every residual cost unperturbed, so ties
and flat directions among residual variables can persist on every draw.

Two marginal laws are used. The first is the grid law $U_{\sigma,M}$. The
second is a \emph{Gaussian-like} finite rational law, constructed in
\cref{sec:qp}, whose Kolmogorov distance to the normal law $N(0,\sigma^2)$ is
at most $2^{-b}$ and which is produced by a bounded rejection sampler using
polynomially many fair bits. The exact normal law is used only as a proof
device; it is not an input model, because its samples are not finite
rational data.

\paragraph{Base-chosen laws.}
In every theorem the law is computed deterministically from the base
instance before sampling: the integer $M$, or the accuracy and support
parameters of the Gaussian-like law, are functions of the base data. The law
therefore depends on the instance. The theorems assert their bounds for this
constructed law. They do not assert the same bounds for every finite law of
half-width $\sigma$, for a coarse grid chosen independently of the instance,
or for an arbitrary approximation of a Gaussian. The reason is the atomic
term in \eqref{eq:model:interval}: the law must be fine enough that atoms on
exceptional sets have small total probability, and ``fine enough'' is
measured against base-only quantities such as the cost of an exact fallback.

\paragraph{Sampling precision.}
Write $b=\log_2M$ for the number of random bits per perturbed coordinate. A
theorem has \emph{polynomial sampling precision} if $b\le\poly(I)$, so the
sampled instance has length $I+n\cdot\poly(I)$. A few results for nonlinear
flows and totally unimodular systems need \emph{parameter-dependent sampling
precision} $b\le f(k)\poly(I)$; this enlarges the sampled instance, and the
stated work bounds account for it.

\paragraph{Small and strong noise.}
In most results $\sigma$ is a free parameter, and the bounds degrade as
$\sigma$ decreases; these are statements about small perturbations of an
arbitrary instance. The random-component results of \cref{sec:int} instead
require noise that dominates the variation of every coordinate derivative.
There $\sigma$ is large compared with the curvature of the instance, and the
results describe instances whose linear part is strongly random rather than
slightly perturbed.

\subsection{Exact smoothed algorithms}\label{sec:model:algorithms}

Computation is measured in bit operations of a deterministic Turing
machine. Oracle calls, where a theorem allows them, are charged at their
stated cost.

\begin{definition}[Exact smoothed algorithm]\label{def:model:algorithm}
Fix a class of base instances, an output format, a perturbation model, and a
rule assigning to each base instance $\Pi$ a finite rational law
$\mathcal D_\Pi$ for $\gamma$. An \emph{exact smoothed algorithm with expected
work $W$} is a deterministic algorithm $\mathcal A$ with the following
properties.
\begin{enumerate}[label=(\roman*)]
\item\label{it:model:every} For every base instance $\Pi$ and every $\gamma$
in the support of $\mathcal D_\Pi$, $\mathcal A(\Pi,\gamma)$ halts and
returns an exact global optimizer and the optimal value of $F_\gamma$ on $X$,
in the stated format.
\item\label{it:model:once} The vector $\gamma$ is drawn once. The algorithm
does not resample, discard, or condition on any event of the draw.
\item\label{it:model:expect} The bit work of computing $\mathcal D_\Pi$,
sampling $\gamma$, and running $\mathcal A(\Pi,\gamma)$ has expectation at
most $W(\Pi)$ over $\gamma\sim\mathcal D_\Pi$.
\end{enumerate}
\end{definition}

The randomness defines the instance to be solved; the algorithm itself is
deterministic given the draw. Property~\ref{it:model:every} includes draws
with tied or positive-dimensional optimal sets and draws on which the fast
route fails. Property~\ref{it:model:once} excludes the common shortcut of
redrawing until a draw is generic, which would change the law of the solved
instance and with it the meaning of the expectation. By Markov's
inequality, an expected bound $W$ implies work at most $tW$ with
probability at least $1-1/t$; it is not a worst-case bound, and on rare
draws the work can be exponential in $I$.

\paragraph{Same-draw fallback.}
Every positive result combines a fast route with an exact
\emph{fallback}: a deterministic algorithm that solves every sampled
instance exactly, with bit work at most $B\cdot\poly(I+b)$, where the factor
$B$ depends only on the base instance and has $\log B\le\poly(I)$. The
fallback runs on the same draw when the fast route has not finished by a
stopping depth $J$ computed from base data. If the fast route can fail only
on an event of probability at most $1/B$, the fallback contributes at most
$\poly(I+b)$ to the expected work. The choices are made in a fixed order:
first the fallback factor $B$, then the thresholds defining exceptional
events, then $J$, and finally the law, for instance the grid size $M$. Only
the polynomial factor of the fallback depends on $b$. This order avoids a
circular dependence between the sampling precision and the precision the
analysis requires.

\paragraph{Expectations of adaptive searches.}
The fast routes refine cells adaptively, and which cells are visited
depends on the draw. Expected counts are therefore never computed as if the
retained cells were independent of the noise. Instead, every visited cell is
charged to a node of a deterministic grid fixed before sampling, and the
expected number of nodes that satisfy a necessary near-optimality condition
is bounded node by node.

\subsection{Output contracts}\label{sec:model:outputs}

Exact output for $F_\gamma$ takes one of the following forms. Each theorem
states which form it returns on ordinary draws and on fallback draws.

\begin{definition}[Output formats]\label{def:model:outputs}
\leavevmode
\begin{enumerate}[label=(\alph*)]
\item\label{it:model:rational} \emph{Rational output}: a feasible
$x^*\in\Q^n$ with $F_\gamma(x^*)=\min_XF_\gamma$, and this value.
\item\label{it:model:implicit} \emph{Implicit patch output}: a set $S$ of
coordinates with fixed rational values $z_S$ (native integer labels and
original continuous bounds), a rational box $Q$ in the remaining
coordinates such that $(z_S,y)\in X$ for every $y\in Q$, and a rational
$g_0>0$ with a verified proof that
$\nabla^2_{yy}F_\gamma(z_S,y)\succeq g_0I$ on $Q$, together with the claim
that the unique minimizer $\hat y$ of $F_\gamma(z_S,\cdot)$ on $Q$ gives a
global optimizer $\hat x=(z_S,\hat y)$ of $F_\gamma$ on $X$. Products of
simplices and order polytopes use a rational relative polytope in place of
$Q$.
\item\label{it:model:algebraic} \emph{Algebraic output}: a real algebraic
number $\theta$, given by a nonzero integer polynomial and an isolating
interval with rational endpoints, and rational polynomial maps that give
every continuous coordinate of one optimizer and the optimal value as
functions of $\theta$, with integer coordinates listed explicitly.
\item\label{it:model:component} \emph{Component output}: a partition of the
unfixed coordinates into components, an algebraic output for each
component, rational values for the remaining coordinates, and the optimal
value written as a rational constant plus the sum of the component values.
\end{enumerate}
\end{definition}

Several distinctions are essential.

\emph{Descriptor versus proof.} An implicit patch output is a compact
descriptor of polynomial length. It denotes a unique point because the
restricted objective is strongly convex on $Q$; that point is a global
optimizer because of a separate \emph{global proof record}, which contains
the pruning trace and the certificate checks. The descriptor alone does not
prove global optimality. The proof record can be large on an individual
draw; its expected size obeys the same bound as the expected work.

\emph{Evaluation.} Given the descriptor and an integer $q\ge0$, a
\emph{$q$-bit evaluation} returns a rational point within distance $2^{-q}$
of $\hat x$ and a rational interval of width at most $2^{-q}$ containing the
optimal value. For patch outputs this costs $\poly_d(I+q)$ bit operations:
convex optimization in the bit model gives value accuracy, and the verified
modulus $g_0$ converts value accuracy into distance. Evaluation does not
decide exact comparisons of the optimal value with a rational threshold,
and it does not identify the exact active set.

\emph{Algebraic sizes.} An algebraic output fixes one optimizer: all its
coordinates are functions of one root $\theta$, so coordinates of different
optimizers cannot be mixed. Its defining polynomial need not be minimal.
Algebraic outputs from fallback draws can have size exponential in $I$, with
polynomial expected contribution. A component output does not include a
defining polynomial of the total value, whose degree can grow exponentially
with the number of components, and it does not provide exact sign or
equality tests for sums of unrelated algebraic numbers. Rational
enclosures of the total value are obtained by refining each component.

\emph{Selection.} Unless a theorem says so, the returned optimizer is the
one the algorithm happens to certify. No canonical selection rule, such as
the lexicographically least or minimum-norm optimizer, is promised.

\subsection{Fixed-parameter and fixed-width bounds}\label{sec:model:parameters}

A bound is \emph{fixed-parameter} in a parameter $\kappa$ if it has the form
$f(\kappa)\cdot\poly(I)$, where the exponent of the polynomial does not
depend on $\kappa$. The low-negative-inertia, core, and integer-dimension
results are of this form, with numerical ratios such as
$\nu\diam(X)/\sigma$ or $L/\sigma$ included in $\kappa$. A bound is
\emph{fixed-width polynomial} if it is polynomial for each fixed value of a
parameter $p$ but its degree grows with $p$, for instance
$[nLw_{\max}/\sigma]^{O(p)}\poly(I)$. The sparse results are of this kind,
and they are not fixed-parameter bounds in the width. The limitation is
genuine for the method used: \cref{thm:lim:width} exhibits instances on
which its retained state count is $(\Omega(n/p))^{p/2}$ on every draw.

A bound that is polynomial in a numerical ratio is polynomial in its
magnitude, which can be exponential in its encoding length. The phrase
\emph{expected polynomial time} is used only when the relevant ratios are
polynomially bounded in $I$.

\subsection{The original objective}\label{sec:model:original}

The theorems optimize the sampled objective $F_\gamma$. An exact optimizer
of $F_\gamma$ need not optimize $F_0$, and nothing is claimed about its
distance from the optimal set of $F_0$. The following elementary bound states
what a sampled solution certifies for the original problem. For compact $X$
and $\gamma\in\R^n$, let
\[
 \omega_X(\gamma)=\max_{x\in X}\gamma^{\mathsf T}x-\min_{x\in X}\gamma^{\mathsf T}x
\]
be the width of $X$ in the direction $\gamma$.

\begin{proposition}[Certified original-objective regret]\label{prop:model:regret}
Let $X$ be compact, let $F_0$ be continuous, and let $y\in X$ and $a\in\R$
satisfy
\[
 a\le\min_XF_\gamma\le F_\gamma(y)\le a+\delta.
\]
Then the interval with endpoints
\[
 a-\max_{x\in X}\gamma^{\mathsf T}x\qquad\text{and}\qquad F_0(y)
\]
contains $\min_XF_0$, and its length is at most $\delta+\omega_X(\gamma)$. In
particular $F_0(y)-\min_XF_0\le\delta+\omega_X(\gamma)$.
\end{proposition}

\begin{proof}
For every $x\in X$, $F_0(x)=F_\gamma(x)-\gamma^{\mathsf T}x\ge
a-\max_X\gamma^{\mathsf T}x$, which gives the lower endpoint. The upper
endpoint is the value of the feasible point $y$. Finally,
\[
 F_0(y)=F_\gamma(y)-\gamma^{\mathsf T}y\le a+\delta-\min_X\gamma^{\mathsf T}x
 =\Bigl(a-\max_X\gamma^{\mathsf T}x\Bigr)+\delta+\omega_X(\gamma).\qedhere
\]
\end{proof}

For an exact sampled optimizer, $\delta=0$ and $a=\min_XF_\gamma$. The width
$\omega_X(\gamma)$ is at most $\sigma\sum_iw_i$ for ambient grid noise on a
mixed box, at most $\sigma\sum_{i\in C}w_i$ ($=k\sigma$ on a unit core box)
for core-only noise, and at most $\sigma\sum_{i=1}^k\omega_X(T_{i\cdot})$ for
aligned noise, where $T_{i\cdot}$ is the $i$th row of $T$. The endpoints are
rational numbers that can be computed exactly: on a box,
$\max_X\gamma^{\mathsf T}x$ is an endpoint calculation; on a polytope it is a
linear program. On a mixed polytope one may maximize over its continuous
relaxation instead, which keeps the lower endpoint valid and replaces
$\omega_X$ by the corresponding width of the relaxation. For Gaussian-like
laws the support is bounded, and the same bound holds with $\sigma$ replaced
by the largest support magnitude.

Choosing $\sigma$ proportional to a target accuracy $\varepsilon$ therefore
turns each theorem into an algorithm that returns, on every draw, a feasible
point and a certificate that its original objective value is within
$\varepsilon$ of $\min_XF_0$. The expected work is then polynomial in
$1/\varepsilon$ with a degree that depends on the structural parameter. We
do not claim that this improves on dedicated approximation algorithms, and
the results do not yield exact optimization of $F_0$: their bounds grow as
$\sigma$ decreases. This is consistent with the worst-case hardness of the
unperturbed problems discussed in \cref{sec:intro} and \cref{sec:lim}.

\subsection{Why a finite law forces an exact fallback}\label{sec:model:atom}

Under a continuous noise law, ties among optimizers and other degenerate
configurations typically have probability zero. Under a finite law they can
have positive probability that no refinement of the search removes. The
following example is the smallest instance of this phenomenon; every
positive result of the paper handles it through its same-draw fallback.

\begin{example}[A tie with probability $1/M$]\label{ex:model:tie}
Let $X=[0,1]^2$, $F_0(x)=2x_1x_2-x_1-x_2$, and let $\gamma_1,\gamma_2$ be
independent with law $U_{\sigma,M}$, where $0<\sigma\le1/4$. The event
$\gamma_1=\gamma_2$ has probability exactly $M\cdot M^{-2}=1/M$.

On this event, put $s=1-\gamma_1\in[3/4,5/4]$. The minimum value of
$F_\gamma$ is $-s$, and
\[
 F_\gamma(x)+s=s(1-x_1-x_2)+2x_1x_2 .
\]
If $x_1+x_2\le1$, both terms are nonnegative, and they vanish together only
at $(1,0)$ and $(0,1)$. If $x_1+x_2>1$, then $x_1x_2>0$, and
$(1-x_1)(1-x_2)\ge0$ gives $x_1+x_2-1\le x_1x_2$, so
$F_\gamma(x)+s\ge(2-s)x_1x_2>0$. Hence $F_\gamma$ has exactly the two
global minimizers $(1,0)$ and $(0,1)$ on every draw of this event.

Consequently no certificate that exhibits a strongly convex patch containing
every global optimizer can succeed on this event, at any refinement depth:
a strictly convex function has at most one minimizer on a convex set. The
point-growth modulus at the optimum is zero, so any tail bound of the form
$\Pr\{g_*<\varepsilon\}\le c\varepsilon+\beta$ for the grid law must have
$\beta\ge1/M$ for all $\varepsilon>0$. An algorithm that is correct on every
draw must therefore handle this event by other means; in this paper that is
the exact fallback, whose expected contribution is at most
$M^{-1}$ times its cost on the event. Choosing $M$ large compared with the
base-only fallback factor $B$ keeps this contribution polynomial. Under any
absolutely continuous law for $(\gamma_1,\gamma_2)$ the tie has probability
zero, which is why continuous-noise analyses can ignore it.
\end{example}

The same example shows why the probability estimates of the later sections
have the form ``continuous bound plus atomic term'', and why the grid size
$M$ is chosen only after the fallback factor and the stopping depth are
known.
